%=============================================================================
% ABSTRACT — UIDT v3.9
% TKT-2026-05-12-manuscript-elite-revision Step 2/4
% All numerical values referenced via macros from preamble_v39.tex
%=============================================================================
\begin{abstract}
\noindent
This manuscript presents the Unified Information-Density Theory (\UIDT{})
v3.9, a constructive theoretical framework introducing a fundamental scalar
field $S(x)$ representing \emph{vacuum information density}. The theory
extends standard Yang--Mills dynamics through non-minimal scalar coupling,
generating a strictly positive spectral gap via vacuum condensate mechanisms.

\medskip\noindent
\textbf{Mathematical Core.}
The canonical parameter set is derived from a self-consistent system of three
coupled equations, anchored by the Extended Functional Renormalization Group
(FRG) and the Banach Fixed-Point Theorem. Within the formal \textsc{uidt}
axiom system, this construction yields a Yang--Mills spectral gap
$\Delta^* = \DeltaGapFull$ \evA{}, a coupling constant
$\kappa = (\Kappa \pm \KappaUnc)$ \evA{}, and the universal kinetic
invariant $\gamma_{\mathrm{kin}} \approx \GammaKin$ \evAminus{}.
Numerical verification achieves closure residuals ${<}10^{-40}$ at
60-digit working precision (\texttt{mpmath}, \texttt{dps\,=\,60}).

\medskip\noindent
\textbf{Four-Pillar Architecture.}
\begin{enumerate}[leftmargin=1.8em,itemsep=2pt]
  \item \textbf{QFT Foundation (Pillar~I).}
        Constructive Banach derivation of $\Delta^*$; RG fixed-point
        identity $5\kappa^2 = 3\lambda_S$ (exact at \texttt{dps\,=\,60});
        benchmark against quenched lattice QCD and QCD sum rules \evB{}.
  \item \textbf{Cosmological Harmony (Pillar~II).}
        Holographic suppression mechanism for the vacuum-energy discrepancy;
        DESI~DR2 calibration of dark-energy equation of state
        $w_0 \approx \wZero$ \evC{};
        Barrow--R\'{e}nyi--Kaniadakis entropy extension;
        Supermassive Dark Seeds (SMDS) model with JWST predictions \evC{}.
  \item \textbf{Laboratory Verification (Pillar~III).}
        Torsion Binding Energy $E_T = \ETorsionfull$ \evC{};
        predicted Casimir anomaly (testable at sub-micron separations);
        geometric correspondence with the bare up-quark mass \evC{}.
  \item \textbf{Photonic Isomorphism (Pillar~IV).}
        Structural mapping of the vacuum lattice to a nonlinear optical
        medium; critical refractive index
        $n_{\mathrm{c}} \equiv \gamma \approx \GammaKin$ \evD{};
        testability via nonlocal metamaterial platforms \evD{}.
\end{enumerate}

\medskip\noindent
\textbf{Epistemic Status.}
All claims are stratified into three layers (Stratum~I: empirical
measurements; Stratum~II: scientific consensus; Stratum~III: \textsc{uidt}
interpretation) and classified by evidence category A--E. Seven explicit
\emph{Kill-Switch} falsification criteria (F1--F7) are provided.
Known limitations---including the phenomenological status of $\gamma$
(Open~Question L4), the absence of independent $E_T$ verification, and
the cosmological calibration uncertainties---are acknowledged throughout
and summarised in Section~\ref{sec:limitations}.

\medskip\noindent
\textbf{Keywords:}
Yang--Mills mass gap; vacuum information density; gamma scaling;
renormalization group; holographic principle; torsion binding energy;
photonic isomorphism; DESI dark energy.
\end{abstract}
